\documentclass[UTF8,a4paper]{ctexart}
\usepackage{amsmath,amssymb,booktabs,geometry}
\geometry{margin=2.4cm}
\title{FSDS $\times$ RAP：Clever Covariate 的效率与经济口径}
\author{}
\date{2026}

\begin{document}
\maketitle

\section{注入位置}
RAP 的选择步使用
\[
H(A,W)=\frac{A-e(W)}{e(W)\,(1-e(W))}.
\]
这里 \(W\) 是当前节点的子节点集合，\(e(W)\) 是稳定分的 softmax。评价器选中的那个动作记 \(A=1\)，于是 \(H=1/e\)。前 12 局只记录 \(H\)，取其 95\% 分位并冻住。其后某一步的 \(H\) 超过该分位，这一步改用稳定分选子节点。节点仍然全部打分，步数不变。

两跳问答的分位是 18.98，Game of 24 的分位是 26.38。漂移后，两跳有 37.5\% 的局至少改了一步，Game of 24 有 45.8\%。

\section{效率}
每局展开的节点数和步数如下。Clever covariate 不减少展开。

\begin{center}
\begin{tabular}{lrrrr}
\toprule
策略 & 漂移前成功率 & 漂移后成功率 & 节点/局 & 步/局 \\
\midrule
两跳，评价器 & 1.00 & 0.25 & 39.0 & 2 \\
两跳，稳定分 & 0.12 & 0.12 & 39.0 & 2 \\
两跳，clever & 0.88 & 0.25 & 39.0 & 2 \\
Game of 24，评价器 & 1.00 & 0.25 & 50.6 & 3 \\
Game of 24，稳定分 & 0.19 & 0.04 & 50.0 & 3 \\
Game of 24，clever & 0.75 & 0.17 & 50.5 & 3 \\
\bottomrule
\end{tabular}
\end{center}

漂移后每节点成功率：两跳评价器与 clever 同为 \(0.0064\)；Game of 24 从 \(0.0049\) 降到 \(0.0033\)。参考期里 clever 已经把两跳从 1.00 打到 0.88，把 Game of 24 从 1.00 打到 0.75，因为第 12 局之后、漂移之前就有步超过了分位。

\section{经济口径}
\[
\mathrm{NetValue}=\Delta\mathrm{Revenue}-\Delta\mathrm{Cost},
\quad
\Delta\mathrm{Revenue}=\Delta\mathrm{Success}\times V\times N,
\quad
\Delta\mathrm{Cost}=0.
\]
\(N=24\) 是漂移后的局数。节点数没有增加，故 \(\Delta\mathrm{Cost}=0\)。\(V\) 只作标价：两跳每局 \(0.001\)，Game of 24 每局 \(0.10\)。没有把 AUC 换成另一笔收入。

\begin{center}
\begin{tabular}{lrrrr}
\toprule
任务 & \(\Delta\)Success & \(\Delta\)Revenue & 未收回的成功率 & 口径 \\
\midrule
两跳 & 0 & 0 & 0.75 & 挽损，未补上 \\
Game of 24 & \(-0.083\) & \(-0.20\) & 0.83 & 挽损，缺口更大 \\
\bottomrule
\end{tabular}
\end{center}

判定规则：clever 的漂移后成功率不超过该任务漂移前评价器的成功率，就记为挽损；超过才记为用增。两题都没有超过漂移前的 1.00。两跳的漂移后成功率与评价器相同，掉下去的 0.75 没有回来。Game of 24 的漂移后成功率比评价器还低 0.083，缺口从 0.75 扩到 0.83。因此这次注入不是用增。

稳定分本身在漂移后更弱（两跳 0.12，Game of 24 0.04）。\(H\) 大时改走稳定分，等于在评价器已经变差的时候再换一条更差的规则。选择步上的异常分可以标出「这一步的动作偏离了先验」，但不能把先验设成一条更弱的启发式。

\end{document}
